\section{Methodology}
\label{sec:methodology}

\subsection{Data Cohort and Feature Processing}
\label{subsec:cohort}

\subsubsection{Study Population and Cohort Selection}
We evaluated adult emergency department (ED) presentations across participating acute care hospital sites. The eligible cohort comprised all adult patient visits ($\ge 18$ years) presenting for acute, unscheduled assessment. We excluded: (i) scheduled follow-up visits or non-emergency elective presentations, (ii) patients certified dead on arrival, and (iii) encounters missing primary outcome adjudication. Detailed center characteristics, inclusion flowcharts, and hospital selection criteria are provided in Supplementary Appendix Section A.

\subsubsection{Predictor Specification and Harmonization}
Predictors were partitioned into two distinct tiers reflecting clinical workflow:
\begin{enumerate}
    \item \textbf{Baseline Triage Predictors ($X_{\text{base}}$):} Baseline administrative and clinical characteristics available at initial triage: age, biological sex, arrival mode, baseline triage urgency category, and chief complaint. Unstructured chief complaints exhibited significant label heterogeneity across sites and were harmonized into standardized diagnostic and syndromic categories via clinical ontology mapping (ICD-10 chapters and SNOMED-CT concepts). Records missing baseline demographic parameters were eliminated during cleaning.
    \item \textbf{Candidate Vital Sign Parameters ($V$):} Point-of-care physiological vitals, including systolic blood pressure (SBP), heart rate (HR), respiratory rate (RR), peripheral oxygen saturation ($\text{SpO}_2$), and body temperature.
\end{enumerate}
Encounters were systematically routed according to feature availability: patients with valid baseline variables were routed to the ED-Atlas baseline prior pipeline, whereas patients satisfying vital acquisition criteria entered the multivariable vital sign integration pipelines. Complete data dictionaries, ontology tables, and handling of unmapped complaints are cataloged in Supplementary Appendix Section B.

---

\subsection{Bayesian Model Architectures}
\label{subsec:bayesian_models}

Let $Y_i \in \{0, 1\}$ denote the binary acute endpoint (e.g., intensive care admission or 30-day mortality) for encounter $i$. Let $X_i \in \mathbb{R}^p$ represent baseline triage covariates and $V_i \in \mathbb{R}^q$ represent physiological vitals.

\subsubsection{Modular Updating (Post-Hoc Integration)}
To retain calibration and transportability established by existing clinical risk models, our modular integration architecture treats the validated ED-Atlas risk assessment as an informative prior. Let $\pi_{\text{Atlas}}(X_i) = P(Y_i = 1 \mid X_i)$ denote the out-of-sample posterior probability generated by the pre-existing ED-Atlas model. This risk is projected onto the unbounded logit scale:
\begin{equation}
    \eta_{0, i} = \text{logit}(\pi_{\text{Atlas}}(X_i)) = \ln\left(\frac{\pi_{\text{Atlas}}(X_i)}{1 - \pi_{\text{Atlas}}(X_i)}\right)
\end{equation}
Vital sign parameters $V_i$ enter through a post-hoc Bayesian evidence-updating stage:
\begin{align}
    \text{logit}(P(Y_i = 1 \mid X_i, V_i)) &= \alpha + \gamma \, \eta_{0, i} + \boldsymbol{\beta}_V^\top g(V_i) \\
    \alpha &\sim \mathcal{N}(0, \sigma^2_\alpha), \quad \gamma \sim \mathcal{N}(1, \sigma^2_\gamma), \quad \boldsymbol{\beta}_V \sim \mathcal{N}(0, \boldsymbol{\Sigma}_V)
\end{align}
where $g(\cdot)$ specifies non-linear parameter transformations (e.g., restricted cubic splines), $\gamma$ acts as a shrinkage parameter centered at $1.0$ to regularize the model toward the historical ED-Atlas calibration, and $\alpha$ accounts for baseline outcome prevalence shifts between cohorts.

\subsubsection{De Novo Bayesian Modeling}
To evaluate whether modular integration constrains predictive capacity, an unconstrained \textit{de novo} joint Bayesian model was formulated to jointly estimate parameters over all baseline predictors and physiological vitals:
\begin{align}
    \text{logit}(P(Y_i = 1 \mid X_i, V_i)) &= \alpha_0 + \boldsymbol{\beta}_X^\top X_i + \boldsymbol{\beta}_V^\top g(V_i) \\
    \alpha_0 &\sim \mathcal{N}(0, 10^2), \quad \boldsymbol{\beta}_X \sim \text{Horseshoe}(\tau, \lambda_X), \quad \boldsymbol{\beta}_V \sim \mathcal{N}(0, \boldsymbol{\Sigma}_V)
\end{align}
Sparse shrinkage priors (regularized horseshoe) were placed on $\boldsymbol{\beta}_X$ to prevent parameter dispersion across high-dimensional, sparse chief complaint indicators.

---

\subsection{Missingness Experiments}
\label{subsec:missingness}

\subsubsection{Handling Informative Missingness}
In acute care, vital sign collection constitutes an intentional clinical action driven by patient acuity and triage protocols. Thus, missing vital parameters violate the Missing Completely at Random (MCAR) assumption and operate under Missing at Random (MAR) conditional on triage covariates, or informative Missing Not at Random (MNAR) mechanisms.

Let $R_{ij} = \mathbb{I}(V_{ij} \text{ is observed})$ denote the observation indicator for patient $i$ and vital $j$. We evaluated four paradigms:
\begin{enumerate}
    \item \textbf{Complete-Case Analysis Baseline:} Serves as a reference benchmark, evaluating performance strictly on the subset with complete vital profiles ($\prod_j R_{ij} = 1$).
    \item \textbf{Domain-Informed Structural Diagnostics:} Clinical practice enforces structured co-observation patterns. Under standardized protocols (e.g., National Early Warning Score, NEWS), automated monitors capture blood pressure, heart rate, and oxygen saturation concurrently ($P(R_{i, \text{HR}} = 1 \mid R_{i, \text{SBP}} = 1) \approx 1$), whereas manual parameters (respiratory rate, body temperature) exhibit protocol-dependent disassociation. Inter-hospital variations in $P(R_{ij} = 1)$ were diagnosed to quantify center-specific collection policies.
    \item \textbf{Missingness Mask as an Acuity Feature:} To leverage informative non-measurement, we parameterized the observation mask directly, weighted under MAR triage assumptions:
    \begin{equation}
        \boldsymbol{\eta}_i = \boldsymbol{\beta}_V^\top \tilde{V}_i + \boldsymbol{\delta}^\top (\mathbf{1} - \mathbf{R}_i) + \boldsymbol{\theta}^\top \left( (\mathbf{1} - \mathbf{R}_i) \odot X_{\text{base}, i} \right)
    \end{equation}
    where $\tilde{V}_i$ denotes median-imputed vital values, and $\boldsymbol{\delta}$ quantifies the independent prognostic hazard associated with the clinical deferral of measurement.
    \item \textbf{Bayesian Imputation Engines:} Fully integrated latent-variable imputation embedded within the Bayesian inference pipeline, treating missing vital parameters as latent parameters sampled jointly conditional on $X_{\text{base}}$.
\end{enumerate}

\subsubsection{Simulation of Sampling-Induced Distribution Shifts}
To evaluate model stability under institutional protocol shifts (e.g., transition from selective physician-ordered testing to universal mandatory nursing triage entry), we conducted synthetic perturbation experiments on the holdout test set. We systematically varied the observation probability:
\begin{equation}
    P(R_{ij}^* = 1 \mid X_i; \kappa) = \text{logit}^{-1}\left( \text{logit}(P(R_{ij} = 1 \mid X_i)) + \kappa \right)
\end{equation}
By sweeping $\kappa \in [-2.5, +2.5]$, we simulated synthetic deployment environments ranging from severe resource under-sampling to universal monitoring. Models trained on historical distributions were evaluated across these shifted test cohorts without retraining to quantify sensitivity to assumption violations.

---

\subsection{Evaluation}
\label{subsec:evaluation}

In accordance with recent methodological guidance on predictive AI for medical decision support (Van Calster et al., \textit{Lancet Digital Health} 2025), model validation was structured across distinct statistical domains (discrimination, calibration, overall performance) and a decision-analytical domain (clinical utility). Improper classification metrics that arbitrarily treat false positives and false negatives as equally costly or depend on arbitrary cutoffs (such as accuracy, Youden index, and F1 score) were excluded from formal performance assessment.

All models were evaluated across the same fixed holdout test cohort. We systematically benchmarked four model tiers:
\begin{enumerate}
    \item $\mathcal{M}_0$ (ED-Atlas Reference): The fixed, established historical baseline model.
    \item $\mathcal{M}_{\text{base}}$ (Retrained Baseline): Re-estimated baseline model using identical $X_{\text{base}}$ predictors on current data to benchmark cohort and temporal drift.
    \item $\mathcal{M}_{\text{modular}}$ (Post-Hoc Bayesian Model): Historical ED-Atlas prior updated with vitals $V$.
    \item $\mathcal{M}_{\text{direct}}$ (\textit{De Novo} Model): Jointly estimated model containing $X_{\text{base}} + V$.
\end{enumerate}

\subsubsection{Evaluation Hypotheses}
We formulated three sequential comparative tests:
\begin{enumerate}
    \item \textbf{Baseline Verification (3 Base Variables Test):} Compare $\mathcal{M}_{\text{base}}$ against $\mathcal{M}_0$. A retrained baseline model must demonstrate parity with the historical ED-Atlas reference to confirm that apparent gains from vitals are not artifacts of baseline cohort drift.
    \item \textbf{Incremental Value of Vital Sign Parameters:} Test whether incorporating vital signs ($\mathcal{M}_{\text{modular}}$ and $\mathcal{M}_{\text{direct}}$) improves discrimination, calibration, and clinical net benefit compared to the baseline triage models ($\mathcal{M}_0$ and $\mathcal{M}_{\text{base}}$).
    \item \textbf{Modular Updating vs. De Novo Architectures:} Evaluate whether post-hoc Bayesian integration ($\mathcal{M}_{\text{modular}}$) matches or exceeds unconstrained \textit{de novo} joint fitting ($\mathcal{M}_{\text{direct}}$). This determines whether healthcare systems can incorporate physiology modularly while preserving the calibrated historical baseline prior.
\end{enumerate}

\subsubsection{Evaluation Metrics}
Uncertainty was quantified using 95\% confidence intervals (CIs) derived via 1,000 stratified bootstrap resamples:
\begin{itemize}
    \item \textbf{Discrimination:} Area under the Receiver Operating Characteristic curve (AUROC / C-statistic), reported as overall micro-averaged AUROC and macro-averaged across primary chief complaint categories. In line with Van Calster et al., rank-based AUROC is retained as the proper statistical measure of discrimination, whereas threshold-conflated measures (such as AUPRC and partial AUROC) are omitted as they mix statistical and decision-analytical domains without following decision-theoretic principles.
    \item \textbf{Calibration:}
    \begin{itemize}
        \item \textit{Observed-to-Expected Ratio ($O:E$):} Quantifies mean calibration (calibration-in-the-large), where an ideal ratio is $1.0$.
        \item \textit{Calibration Intercept ($\alpha_{\text{cal}}$) and Slope ($\beta_{\text{cal}}$):} Estimated via logistic recalibration ($\text{logit}(Y) = \alpha_{\text{cal}} + \beta_{\text{cal}}\,\text{logit}(\hat{p})$). $\alpha_{\text{cal}} = 0$ indicates accurate overall calibration, while $\beta_{\text{cal}} = 1$ indicates correct spread; $\beta_{\text{cal}} < 1$ signals overfitting.
    \end{itemize}
    \item \textbf{Overall Performance (Strictly Proper Scoring Rules):}
    \begin{itemize}
        \item Out-of-sample log-likelihood and Brier score ($N^{-1}\sum_{i=1}^N (\hat{p}_i - Y_i)^2$). Both represent strictly proper scoring rules combining discrimination and calibration.
        \item Nagelkerke's pseudo-$R^2$ and scaled Brier score (Index of Prediction Accuracy) to evaluate the proportion of explained variation relative to null models.
    \end{itemize}
    \item \textbf{Clinical Utility:} Evaluated using decision-analytic metrics that explicitly incorporate the trade-off between false positives and false negatives across decision thresholds $p_t$:
    \begin{itemize}
        \item \textit{Net Benefit (NB):}
        \begin{equation}
            \text{NB}(p_t) = \frac{\text{TP}}{N} - \frac{\text{FP}}{N}\left(\frac{p_t}{1 - p_t}\right)
        \end{equation}
        \item \textit{Standardized Net Benefit (sNB):} Computed as $\text{sNB}(p_t) = \text{NB}(p_t) / \text{prevalence}$, scaled to a theoretical maximum of 1.
        \item \textit{Expected Cost (EC):} Evaluated across normalized false negative costs $w_{\text{FN}} \in (0, 1)$ with $w_{\text{FP}} = 1 - w_{\text{FN}}$, tracking the threshold that minimizes clinical expected loss.
    \end{itemize}
\end{itemize}

\subsubsection{Graphical Diagnostics}
In accordance with recommended reporting standards, models were visually evaluated via four graphical diagnostics:
\begin{enumerate}
    \item \textbf{Reliability Diagrams (Calibration Plots):} Moderate calibration was assessed using locally weighted scatterplot smoothing (loess) across predicted probabilities with 95\% confidence bands, plotted against the 45-degree line of perfect calibration.
    \item \textbf{Risk Distribution Profiles:} Overlaid violin and kernel density plots of estimated probabilities separated by outcome ($Y=0$ vs. $Y=1$) to characterize risk separation.
    \item \textbf{Decision Curves:} Net Benefit and Standardized Net Benefit plotted across a clinically plausible range of emergency decision thresholds ($p_t \in [0.01, 0.30]$), comparing each model candidate against standard reference strategies (``treat all'' and ``treat none'').
    \item \textbf{Expected Cost Curves:} Expected clinical cost plotted as a function of the normalized misclassification cost ratio to confirm cost minimization across operational environments.
\end{enumerate}